\section{Commitment Bases for Two-Phase Locking and Multiversion Timestamp Ordering}
\label{app:cc-bases}

This appendix gives the transaction specification of
Sections~\ref{sec:ambiguity-story} and~\ref{sec:systems} in full, and the
commitment bases of 2PL and MVTO.

\paragraph{Specification.}
An input is a finite set of transactions, each listing the keys it reads and
the keys it writes; inputs are ordered by inclusion.  An output has three
parts: the committed transactions, in a serial order; the aborted
transactions; and the read results reported so far, each pairing a read of a
key by a committed transaction with the write whose value it returned, or with
the initial version of the key.  One output \outrel another when it extends all
three: it keeps every committed and every aborted transaction, keeps the serial
order of the committed ones, and keeps every reported read result.  A row is
permitted when every transaction of the input is committed or aborted, and not
both, and each reported read of a key $x$ by a transaction $T$ returns the write
of $x$ by the last transaction before $T$ in the serial order that writes $x$,
or the initial version of $x$ if no earlier transaction writes it.  The figures
of Sections~\ref{sec:ambiguity-story}--\ref{sec:systems} show only the serial
order, for transactions without reads.

We write an output as $(L,X,\rho)$: $X$ is the set of aborted transactions, $L$
is the serial order, a total order on the set $\mathrm{dom}\,L$ of committed
ones, and $\rho$ is the set of reported read results.  Then
\[
  (L,X,\rho)\OutputOrd(L',X',\rho')
  \;\Longleftrightarrow\;
  X\subseteq X',\ \ \mathrm{dom}\,L\subseteq\mathrm{dom}\,L',\ \
  L'|_{\mathrm{dom}\,L}=L,\ \ \rho\subseteq\rho'.
\]
The rows $(\{A,B\},([A,B],\varnothing,\varnothing))$ and
$(\{A,B\},([B,A],\varnothing,\varnothing))$, corresponding to
Figure~\ref{fig:ambiguity-frames}(f), have no common \outnoun, since it would
order $A$ before $B$ and $B$ before $A$; so the specification is not monotone.
The \emph{abort-free} specification has the same rows with $X=\varnothing$.

\paragraph{Commitments.}
Each commitment below is a condition on the serial order $L$ and abort set $X$
of a row.  Applied at a state $(i,R)$, it retains the rows of $R$ that satisfy
the condition, and is defined when some row of $R$ at $i$ satisfies it and some
row of $R$ fails it.
\begin{itemize}
  \item $\mathsf{order}(A,B)$, for distinct $A,B\in i$: $A\in X$, or $B\in X$,
        or $A<_L B$.
  \item $\mathsf{commit}(A)$, for $A\in i$: $A\notin X$.\quad
        $\mathsf{abort}(A)$, for $A\in i$: $A\in X$.
  \item $\mathsf{stamp}(A)$, for $A\in i\setminus Z$, given $Z\subseteq i$, the
        transactions stamped before $A$: at row $(k,(L,X,\rho))$, for every
        $z\in Z$, $z<_L A$ or one of $z,A$ is in $X$; and for every
        $D\in k\setminus(Z\cup\{A\})$, $A<_L D$ or one of $A,D$ is in $X$.
\end{itemize}
Each is a commitment (Definition~\ref{def:commitment}).  Two permitted outputs
at one input, one \outrelpart the other, have the same serial order and abort
set: if $(L,X,\rho)\OutputOrd(L',X',\rho')$ then $X\subseteq X'$ and
$\mathrm{dom}\,L\subseteq\mathrm{dom}\,L'$, and both pairs partition the input,
so $X=X'$ and $\mathrm{dom}\,L=\mathrm{dom}\,L'$, hence
$L=L'|_{\mathrm{dom}\,L}=L'$.  Each condition depends only on $L$ and $X$, so it
holds at both or at neither; a removed row is therefore not \outrelpast by a
retained row at the same input, which is
Equation~\ref{eq:commitment-excludes}.  The bases are
$\mathcal B_{\mathrm{2PL}}=\{\mathsf{order},\mathsf{commit},\mathsf{abort}\}$
and
$\mathcal B_{\mathrm{MVTO}}=\{\mathsf{stamp},\mathsf{commit},\mathsf{abort}\}$,
ranging over all transactions.  A scheduler applies them within the sequence
$\gamma$ of an output event $\Out{\gamma}{o}$
(Definition~\ref{def:system-output-event}); read results enter only through
the emitted output $o$.  The protocol steps below say which commitments
$\gamma$ contains.

\paragraph{2PL.}
When $B$ waits for a lock that $A$ holds in a conflicting mode, the scheduler
applies $\mathsf{order}(A,B)$; on rows without aborts it retains exactly those
with $A$ before $B$ (Figure~\ref{fig:cc-commitments}(h)).  When $A$ commits, the
scheduler applies $\mathsf{commit}(A)$.  A deadlock between $A$ and $B$ means
both $\mathsf{order}(A,B)$ and $\mathsf{order}(B,A)$ have been applied; the
rows left, at $i$ and every larger input, are exactly those with $A\in X$ or
$B\in X$, so every later output aborts $A$ or $B$, and resolving the deadlock is
applying $\mathsf{abort}(A)$ or $\mathsf{abort}(B)$.  Under the abort-free
specification, once $\mathsf{order}(A,B)$ is applied no remaining row
satisfies $\mathsf{order}(B,A)$, now or later, so it is never defined: $A$, and
hence $B$, waits forever.  Deadlock handling thus corresponds to the choice
between the two specifications.

\paragraph{MVTO.}
On arrival $A$ receives a timestamp above all issued ones, and the scheduler
applies $\mathsf{stamp}(A)$, with $Z$ the transactions stamped before $A$; on
rows without aborts, $\mathsf{stamp}(B)$ with $Z=\varnothing$ retains exactly
those with $B$ first (Figure~\ref{fig:cc-commitments}(i)).  Reads add no
commitment.  The scheduler applies $\mathsf{abort}(A)$ when $A$ writes a version
that would directly follow a committed version already read by a transaction
stamped after $A$, and otherwise applies $\mathsf{commit}(A)$ when $A$ commits.
The specification itself requires this abort.  Suppose $A$ writes $x$, $B$ was
stamped after $A$, and an output reports that $B$ read the initial version of
$x$, or the write of $x$ by a transaction $W$ stamped before $A$.  $B$ is
committed, since it has a reported read.  If $A$ were committed too,
$\mathsf{stamp}(B)$ would put $A$ before $B$ in $L$, and $\mathsf{stamp}(A)$
would put $W$ before $A$; then $A$ writes $x$ after the version $B$ read and
before $B$, so the reported read would not be permitted.  So any output
reporting $B$'s read aborts $A$, and so does every row left once it is
emitted.  The writer must be stamped before $A$: at input $\{A,W,B\}$ with
$A$, $W$, $B$ stamped in that order and $A$, $W$ writing $x$, the row that
commits all three in stamp order and reports that $B$ read $W$'s write is
permitted and satisfies every stamp.  The requirement comes from the reported
read: at input $\{W,A,B\}$, stamped in that order, after the three stamps a
row that commits $A$ and reports no read for $B$ remains.

\paragraph{Comparison.}
2PL commits relative orders, only between conflicting transactions; MVTO
commits each transaction's position against all present and future ones on
arrival, and pays with aborts when the reads it reports do not match that
position.  Every output decides each transaction of its input, so fixing only
$A$'s status or its relative position is a commitment, never itself an output.
